%=====================================================================
\chapter{Exercises and Answer Sketches}\label{app:exercises}
%=====================================================================

Sketches for the \emph{Checkpoint} questions in each chapter. They are sketches:
enough to tell whether you were on the right track, not model answers. The
\emph{Review Questions} are deliberately open and have no sketches --- they are
seminar and examination material, and several have no settled answer.

\section{Part I --- Foundations of Inquiry}

\subsection*{\chref{contribution} --- What Counts as a Contribution}
\begin{tightnum}
  \item Not as stated: it is an activity, with no claim and no falsifier. The
    addition that fixes it is a comparison plus a condition --- for example,
    that a model using LMS engagement data predicts dropout better than
    prior-CGPA alone, which could come out either way.
  \item Any case where the field assumes an effect: a well-powered study
    finding that a widely recommended practice does not reduce defects is
    knowledge, and changes what people should do.
  \item An ablation protects against attributing a gain to the wrong component;
    a threats section protects against attributing an observed difference to
    the treatment when something else produced it. Both isolate the cause of
    an effect, in different kinds of study.
  \item Replicate it under a condition the original did not test; test the
    mechanism it assumed but did not verify; or extend it to a population or
    setting where its assumptions plausibly fail.
\end{tightnum}

\subsection*{\chref{philosophy} --- Philosophy of Science}
\begin{tightnum}
  \item Post-positivism. A confidence interval concedes that measurement is
    imperfect and the conclusion provisional, which positivism's verificationism
    does not.
  \item A benchmark disagreeing with a published result may refute the claim
    --- or your JVM version, your hardware, your workload configuration, or an
    assumption the original left implicit. Rejecting the claim requires ruling
    those out.
  \item Abduction generates; deduction and induction test. It is rarely
    described because methods chapters are written retrospectively, when the
    conjecture already looks obvious.
  \item Under interpretivism there is no view from nowhere, so ``objective''
    cannot mean researcher-free. Offer instead confirmability: an audit trail
    showing findings are traceable to data, plus explicit reflexivity about
    your position.
\end{tightnum}

\subsection*{\chref{questions} --- Research Questions and Constructs}
\begin{tightnum}
  \item P: teams deploying web services. I: containerised deployment.
    C: virtual-machine or bare-metal deployment. O: deployment lead time and
    failure rate. C: small teams without dedicated operations staff.
  \item Operationalisation: quality $=$ comment density. Objection: comment
    density measures documentation habit, and can rise when code is unclear ---
    so it may correlate \emph{negatively} with the construct.
  \item Because the direction was one of two available choices. Choosing after
    seeing the data means both directions were live, so the effective test was
    two-tailed while the reported p-value is one-tailed.
  \item Any exploratory question --- for instance, how reviewers decide what to
    comment on. You cannot predict relations among categories you have not yet
    discovered.
\end{tightnum}

\subsection*{\chref{publication} --- The Publication System}
\begin{tightnum}
  \item Only the post-approval paper counts. They need one more X-category
    paper, or a single W-category paper.
  \item Citation distributions within a journal are heavily skewed: most papers
    are cited far below the journal's mean, so the journal's average says
    little about any individual article.
  \item Better: stronger review, and the audience is there. It may still not
    help because \S15 is written in terms of journals, so a conference paper
    does not satisfy the requirement.
  \item Check HJRS category, DOAJ listing, and Think.\ Check.\ Submit. A
    72-hour review promise and an unrecognised ``impact factor'' would almost
    certainly fail all three.
\end{tightnum}

\section{Part II --- Finding and Synthesising Knowledge}

\subsection*{\chref{slr} --- Systematic Literature Review}
\begin{tightnum}
  \item So that inclusion criteria cannot be adjusted after seeing which papers
    support your preferred conclusion --- the same selection bias
    pre-registration prevents.
  \item Add the Comparison element; restrict fields to title-abstract-keyword
    rather than full text; narrow the date range. Each risks losing relevant
    studies, and the first risks least.
  \item The breakdown of full-text exclusions by reason. Without it a reader
    cannot tell whether criteria were applied consistently or invented during
    screening.
  \item Because studies differ in power. Four small studies finding nothing are
    fully compatible with a real effect they were too small to detect, and nine
    small studies finding one are compatible with publication bias.
\end{tightnum}

\subsection*{\chref{greylit} --- Snowballing and Grey Literature}
\begin{tightnum}
  \item One group's citation neighbourhood is self-reinforcing: they cite each
    other and the same predecessors, so iteration converges without ever
    reaching adjacent communities using different terminology.
  \item Any practice-led topic --- container orchestration, site reliability
    engineering, MLOps --- where industrial practice was documented in
    engineering blogs years before academic study.
  \item Include it as evidence of vendor claims, never as evidence of
    performance, and label it that way. Its objectivity score under the grey
    literature criteria is the lowest possible.
  \item Where a negative result is commercially or politically convenient ---
    for instance, an evaluation finding that a mandated tool has no effect may
    be easier to publish than one finding it helps.
\end{tightnum}

\subsection*{\chref{bibliometrics} --- Reference Management and Bibliometrics}
\begin{tightnum}
  \item Coupling depends on the references a paper \emph{makes}, which exist at
    publication. Co-citation depends on later papers citing it, which takes
    years.
  \item PRISMA counts fall out of the collection structure rather than a
    separate spreadsheet; and exported \texttt{.bib} can be version-controlled
    with the thesis, so the bibliography is reproducible.
  \item Evidence that somebody read papers in each cluster and confirmed the
    label --- ideally with the cluster's most locally cited papers named.
  \item Both index journals from some regions and languages far more completely
    than others, so apparent output differences partly measure indexing
    coverage rather than research activity.
\end{tightnum}

\subsection*{\chref{appraisal} --- Reading and Appraising}
\begin{tightnum}
  \item 900 first passes, roughly 50 second passes, and third passes only for
    the three or four papers your delta is measured against. Second-passing
    during screening costs weeks and changes no screening decision.
  \item Practical, not moral: a criticism that does not survive the strongest
    reading was a misreading, and offering it in a seminar or referee report
    costs you credibility.
  \item ``Result: $p = 0.03$, effect size not reported.'' It matters because
    that study cannot enter a meta-analysis, which is itself worth recording as
    a finding about the field.
  \item Check your own setup; check versions and seeds; re-read the paper for an
    unstated parameter; write to the authors. Only then consider that the
    result may be wrong.
\end{tightnum}

\section{Part III --- Research Designs}

\subsection*{\chref{designscience} --- Design Science Research}
\begin{tightnum}
  \item Formative-naturalistic at best, and arguably only demonstration. The
    objection: no comparison, no stated criteria, and no result that would have
    counted as failure.
  \item Boundary case: applying a standard CNN to a new local crop dataset. It
    is exaptation if local conditions broke the standard approach in a specific
    way you characterised and addressed; routine design if it worked as
    expected.
  \item Otherwise the criteria are selected to be ones the artefact passes ---
    the same bias as choosing an analysis after seeing the data.
  \item The method is the contribution. Structure the thesis so the method has
    its own chapter, with the instantiation presented afterwards as evidence
    that it is realisable.
\end{tightnum}

\subsection*{\chref{experiments} --- Controlled Experiments}
\begin{tightnum}
  \item It makes assignment independent of every variable, including unmeasured
    ones. Statistical adjustment can only handle variables you measured.
  \item Within-subjects, so each developer sees all three techniques. The threat
    is carryover --- and learning across the three sessions --- which
    counterbalancing by Latin square addresses.
  \item A technique-by-artefact interaction: a technique uniquely suited to one
    artefact. Detecting it requires a full factorial design.
  \item Because a rule chosen after seeing which direction it moves the result
    is a researcher degree of freedom, however reasonable the rule is in the
    abstract.
\end{tightnum}

\subsection*{\chref{quasiexp} --- Quasi-Experiments and Mining}
\begin{tightnum}
  \item Because bias does not shrink with $n$. A larger sample estimates the
    confounded quantity more precisely.
  \item Difference-in-differences removes all \emph{time-invariant} unmeasured
    differences between groups; matching removes only measured ones. It fails
    to remove anything that changed at the same time as treatment.
  \item Only some fixes are linked to issues, and linking correlates with
    project discipline; and fixes for defects never reported as issues are
    absent entirely.
  \item Only when no before-and-after data exists --- a purely cross-sectional
    comparison --- and even then the no-unobserved-confounding assumption must
    be stated as the study's principal limitation.
\end{tightnum}

\subsection*{\chref{casestudy} --- Case Study Research}
\begin{tightnum}
  \item Because separating them would destroy the phenomenon. What is given up
    is control, and therefore any strong causal warrant.
  \item Four cases cannot support a statistical claim about a country. Rewrite
    as an analytic generalisation: a named mechanism, its scope conditions, and
    a prediction about where it should and should not appear.
  \item Because it predicts a \emph{difference} in advance and can therefore be
    wrong. Literal replication can only accumulate confirmations.
  \item Either the sources measure different things, or people's accounts
    differ from their behaviour, or one source is unreliable. Distinguish by
    triangulating a third source and by asking participants about the
    discrepancy directly.
\end{tightnum}

\subsection*{\chref{survey} --- Survey Research}
\begin{tightnum}
  \item Population: all Pakistani software professionals. Frame: whatever list
    you can reach, probably online communities. Achieved: who answered. The
    largest error is coverage --- the gap between population and frame.
  \item If the 30\% who declined differ systematically on the measured variable
    --- for instance, if the topic is uncomfortable for exactly those people
    --- the bias can exceed that of a low-rate but indifferent nonresponse.
  \item Split it: ``Compared with the previous pipeline, how has build duration
    changed?'' and ``\ldots how has build reliability changed?'', each with a
    response scale that permits ``worse''.
  \item Several items measuring the one construct, so the summed scale can be
    treated as approximately interval.
\end{tightnum}

\subsection*{\chref{qualitative} --- Qualitative Methods and Coding}
\begin{tightnum}
  \item Because a specific incident is recalled, whereas a general question
    elicits the official or idealised account --- which is frequently not what
    happens.
  \item New codes per interview, showing the curve flattening, plus a stopping
    rule stated in advance. Or an honest statement that access, not saturation,
    determined the number.
  \item A category error where codes are interpretive constructions under a
    constructivist position; a weakness where codes are descriptive categories
    meant to be applied consistently.
  \item ``We constructed three themes from the data.'' Then say who coded, and
    how disagreements were resolved.
\end{tightnum}

\subsection*{\chref{mixed} --- Mixed Methods}
\begin{tightnum}
  \item Because nothing integrates: neither strand informs the other's design,
    sampling or interpretation, so the combined inference is no more than the
    two separate ones.
  \item Explanatory sequential. The quantitative result determines who is
    interviewed --- typically extreme and contrasting cases from among the
    respondents.
  \item That the qualitative strand came first and dominates, and the
    quantitative strand is secondary --- so the reader knows which strand
    carries the claim and can judge its evidence accordingly.
  \item Different constructs, different populations, one strand wrong, or two
    real mechanisms in opposite directions. A subgroup analysis of the
    quantitative strand along the dimension the qualitative strand identified
    usually discriminates.
\end{tightnum}

\subsection*{\chref{benchmarking} --- Benchmarking and Simulation}
\begin{tightnum}
  \item Because the answer changes depending on which system you place in the
    numerator. Compute it both ways: a valid summary gives consistent answers.
  \item Verification asks whether the model was implemented correctly;
    validation asks whether it corresponds to reality. An unvalidated
    simulation may pass verification completely.
  \item Because JIT decisions, memory layout and address randomisation vary
    between processes, and iterations within one process share them --- so
    within-process repetition underestimates true variability.
  \item Report both, prominently. You have discovered evidence that the suite
    is no longer representative for this class of workload, which is a
    contribution in itself.
\end{tightnum}

\section{Part IV --- Measurement, Data and Analysis}

\subsection*{\chref{measurement} --- Measurement and Validity}
\begin{tightnum}
  \item It establishes that the items correlate. It does not establish
    unidimensionality, and with 20 items a high value is close to guaranteed
    regardless of quality.
  \item Because reliability is consistency and validity is correctness. Commit
    count is highly reliable and measures commit style rather than productivity.
  \item A single item is ordinal, so the mean assumes equal spacing between
    response options. Report the median and the full distribution.
  \item Any metric used in performance review --- commits, story points, review
    comments. Detect it by comparing periods before and after the measure
    became visible to participants.
\end{tightnum}

\subsection*{\chref{sampling} --- Sampling and Power}
\begin{tightnum}
  \item It is a deterministic function of the p-value, so it conveys nothing
    beyond it. Report the confidence interval on the effect, or an equivalence
    test.
  \item Under low power only unusually large sample effects reach significance,
    so published effects are systematically over-estimates --- and subsequent
    studies powered from them are underpowered too.
  \item Because small-sample effect estimates are extremely noisy; if the pilot
    over-estimated, you power for an effect that may not exist.
  \item Upwards. Unreliable measurement attenuates the observed effect, so a
    larger sample is needed to detect the same true effect.
\end{tightnum}

\subsection*{\chref{inference} --- Statistical Inference}
\begin{tightnum}
  \item The probability of data at least this extreme \emph{if the null
    hypothesis and every other model assumption held}. The final clause is the
    one usually dropped.
  \item The first is strong evidence the effect is negligible; the second says
    the study was uninformative. Identical p-value verdicts, entirely different
    findings.
  \item Its power runs opposite to your need: it fails to detect non-normality
    in small samples, where it matters, and flags trivial departures in large
    ones, where robustness makes it irrelevant.
  \item Non-significance is not evidence of equivalence. Run a TOST equivalence
    test against a justified bound --- your smallest effect size of interest.
\end{tightnum}

\subsection*{\chref{regression} --- Regression and Modelling}
\begin{tightnum}
  \item Because it is defined as the effect holding the other terms fixed;
    change the set of other terms and you are conditioning on something
    different.
  \item Negative binomial. Poisson would give standard errors far too narrow,
    producing over-confident intervals and p-values.
  \item Whenever the treatment effect plausibly differs between units --- which
    for human participants is nearly always. Omitting it inflates significance
    for the fixed effect.
  \item Because inference requires that the reported standard errors account for
    how the model was chosen, and stepwise selection is an unaccounted analysis
    of the data. Prediction is judged on held-out performance, which is unaffected.
\end{tightnum}

\subsection*{\chref{sem} --- Structural Equation Modelling}
\begin{tightnum}
  \item Random measurement error in a predictor attenuates its estimated
    relationship with the outcome. SEM models the indicators' error explicitly
    and estimates the relationship between the error-free latent constructs.
  \item IT infrastructure maturity, composed of network capacity, staff
    expertise and budget. These need not correlate, so low inter-item
    correlation would wrongly appear as poor reliability.
  \item It suggests redundant items --- several items asking essentially the
    same question --- which inflates reliability without adding measurement
    breadth.
  \item Merge them into a single construct and respecify, disclosing the change
    as data-driven; or redesign the items so the constructs are empirically
    distinguishable and collect again.
\end{tightnum}

\subsection*{\chref{mleval} --- Evaluating Machine Learning}
\begin{tightnum}
  \item The scaler's parameters incorporate test-set statistics, so the model
    sees information unavailable at deployment. The effect is usually small but
    is free to avoid with a pipeline object.
  \item A temporal split, and ideally a grouped one by developer. Random
    $k$-fold would train on the future and let the model memorise individuals.
  \item ROC-AUC is dominated by the abundant true negatives and looks optimistic.
    Report PR-AUC, noting that its baseline is the positive rate, plus
    calibration.
  \item Multiplicity is uncorrected and normality across heterogeneous datasets
    is untenable. Use the Friedman test on ranks, then a post-hoc test, reported
    as a critical difference diagram.
\end{tightnum}

\subsection*{\chref{qda} --- Qualitative Data Analysis}
\begin{tightnum}
  \item Because independent analysts are not expected to produce identical
    themes under an interpretivist position, so replication cannot be the
    standard. Transparency lets a reader judge the route from data to claim.
  \item It invites a prevalence reading of a purposively sampled corpus.
    Legitimate: reporting how many of your participants expressed the theme, as
    a description of your sample only.
  \item Theme: automation as licence to disengage. Interpretation: the
    placement of automated signals in the review interface reframes review from
    evaluation to confirmation --- which predicts that moving them would change
    review depth.
  \item Revise when they identify a factual error or a context you misread. Do
    not revise merely because they dislike an interpretation their own frame
    would not produce --- but report the disagreement.
\end{tightnum}

\subsection*{\chref{validity} --- Threats to Validity}
\begin{tightnum}
  \item Because if there is no relationship in the data, or it is an artefact
    of the analysis, generalising it is meaningless.
  \item Mitigated: randomised assignment removes selection. Measured: attrition
    quantified and compared across conditions. Acknowledged: student subjects,
    with a stated direction of likely bias.
  \item Because it states neither direction nor magnitude, so a reader cannot
    adjust their belief. Replace with the direction and, where possible, a
    bound.
  \item Moving from a laboratory task to a field study in a real company. Right
    when the effect's existence is already established and the open question is
    whether it survives realistic conditions.
\end{tightnum}

\section{Part V --- Rigour, Reproducibility and Ethics}

\subsection*{\chref{reproducibility} --- Reproducibility}
\begin{tightnum}
  \item Repeatability: same team, same setup. Reproducibility: different team,
    your artefacts. Replicability: different team, their own artefacts.
    Re-running your own scripts establishes repeatability only.
  \item Because several libraries hold independent random state, GPU kernels
    may reduce non-deterministically, parallel execution merges in completion
    order, and hash randomisation varies. Report variability across seeds.
  \item Metadata, instruments, protocol, analysis code, and a synthetic or
    aggregated dataset, plus a documented access route. FAIR requires findable
    and accessible metadata, not open data.
  \item Repositories are deleted, renamed, made private and force-pushed, and
    the URL does not identify a version. Archive a release and cite the DOI.
\end{tightnum}

\subsection*{\chref{prereg} --- Pre-registration and QRPs}
\begin{tightnum}
  \item Because the test consisted in the prediction preceding the data. A
    hypothesis fitted to the data it is then tested on has not been tested,
    whatever its truth value.
  \item Because the decision to continue depends on the result so far. Under the
    null, repeated looking guarantees crossing the threshold eventually.
  \item It depends on whether the switch was prompted by the outcome. Prompted
    by a residual plot examined before unblinding: report as planned. Prompted
    by the p-value: relabel as exploratory.
  \item Because publication is already secured, so a significant result confers
    no advantage. The incentive that drives the practice is removed rather than
    prohibited.
\end{tightnum}

\subsection*{\chref{ethics} --- Research Ethics}
\begin{tightnum}
  \item Because voluntariness fails: a student cannot freely decline a request
    from someone who assesses them, whatever they sign.
  \item Commit timing patterns, the specific set of repositories contributed to,
    and coding style. Time zone and commit message idiom are others.
  \item Analysing individuals' commit timestamps to infer religious observance,
    health, or undisclosed employment. Public availability does not make the
    inference an expected use of the data.
  \item Because participants gave time to a study that could not answer its
    question, which is a harm with no offsetting benefit.
\end{tightnum}

\subsection*{\chref{integrity} --- Academic Integrity}
\begin{tightnum}
  \item No. The overall figure is within 19\%, but 9\% from one source breaches
    the ``less than 5\% from a single source'' rule in \S32(III).
  \item Because \S32(I)(A)(ix) names it explicitly: changing words while copying
    the sentence structure. The structure is the author's intellectual work.
  \item An extension adds substantial new material --- typically 30\% or more
    --- discloses the earlier version and cites it. Duplicate publication
    presents the same study as new.
  \item \S15 requires the scholar to be first author for the paper to count.
    Raise it by reference to the regulation rather than as a personal dispute,
    ideally before the work starts.
\end{tightnum}

\subsection*{\chref{genai} --- Generative AI}
\begin{tightnum}
  \item \S32(I)(C)(vi) lists using machine-generated text as plagiarism;
    \S32(IV) permits AI to understand basic phenomena but not to produce
    insights, analyse or interpret. Reading: AI may help you understand; it may
    not produce your contribution or your submitted text.
  \item Because \S32(I)(A)(viii) makes giving incorrect information about the
    source of a quotation plagiarism. A non-existent source is the limiting case.
  \item Because it produces a probability with no matched source to inspect, and
    it flags human writing --- disproportionately from non-native English
    writers. Dated drafts, reading notes and the ability to explain your own
    argument are what answer an accusation.
  \item No. Your thesis is examined under IUB regulations, which are stricter
    here. Follow the stricter rule; disclose in both places.
\end{tightnum}

\section{Part VI --- Producing and Defending}

\subsection*{\chref{synopsis} --- The Synopsis}
\begin{tightnum}
  \item Problem: what is wrong in the world. Rationale: why study it now.
    Significance: who benefits. Novelty: what is new relative to the
    literature. Four different questions.
  \item Because nobody can check whether you explored. Replace with a verb that
    yields a checkable outcome --- measure, compare, design and evaluate.
  \item PhD: end of the sixth semester, failing which the topic is not approved.
    MS: within eight weeks of the start of the third semester, with the title
    fixed by that semester's end.
  \item State the smallest effect you powered for, and that failing to find an
    effect of that size is itself a finding practitioners should act on.
\end{tightnum}

\subsection*{\chref{writing} --- Writing Up}
\begin{tightnum}
  \item Because the thesis carries the through-line, and papers cut from it
    inherit it. Theses assembled from separate papers usually have no argument
    connecting them, and examiners see the seams.
  \item ``\ldots which shows that structured guidance improves reviewer
    attention.'' The mechanism claim belongs in Discussion.
  \item It exposes structure independent of prose. Reading first sentences in
    sequence reveals whether the argument follows, which fluent writing can
    otherwise disguise.
  \item \S14.2(vi), formatting to the IUB guidelines. A non-conforming thesis
    is returned before the research is read.
\end{tightnum}

\subsection*{\chref{publishing} --- Publishing}
\begin{tightnum}
  \item Because it measures whether your readers are there. A venue publishing
    nothing you cite is unlikely to reach the community your work addresses.
  \item To tell the editor why this paper belongs in this venue. Leave out the
    abstract summary, claims of importance, and general motivation.
  \item Correct the reviewer \emph{and} clarify the text, then say you have done
    both. A competent reader misreading it means other readers will too.
  \item It does not count towards \S15, because it was published before synopsis
    approval. The synopsis should have been approved first.
\end{tightnum}

\subsection*{\chref{defending} --- Presenting and Defending}
\begin{tightnum}
  \item Breadth across your specialisation, from the courses you took --- not
    depth in your own topic. Thesis reading is far too narrow.
  \item Because you may be interrupted before reaching the end. The examiners
    must know what you claim early enough to question it.
  \item Concede it plainly, state precisely what it does and does not affect,
    and say what would resolve it. Bluffing is always detected.
  \item PhD: 180 days for the re-defence after corrections, or the candidate is
    considered failed. MS: six months for resubmission from DASR's intimation.
\end{tightnum}

\section{The Synoptic Course Project}

The seven portfolio pieces (Appendix~\ref{app:schedule}) are one continuous
piece of work, not seven exercises. Assembled, they are a defensible synopsis.

\begin{center}
\small
\begin{tabular}{@{}L{1.0cm}L{4.4cm}L{4.2cm}L{4.0cm}@{}}
\toprule
 & \textbf{Piece} & \textbf{Templates} & \textbf{Chapters} \\
\midrule
P1 & Topic statement       & C1, C2      & 1, 3 \\
P2 & SLR protocol          & C3--C6      & 5, 6, 7 \\
P3 & Design choice memo    & C7          & 9--16, 18 \\
P4 & Draft instrument      & C8--C10     & 13, 14, 17, 27 \\
P5 & Analysis plan         & C11         & 19--24 \\
P6 & Pre-registration      & C12, C13    & 25, 26 \\
P7 & Synopsis draft        & C14, C15, C17 & 29, 30, 31 \\
\bottomrule
\end{tabular}
\end{center}

\begin{conceptbox}[Self-assessment before submission]
Score each piece 0--2 and total. Below 10 means the synopsis is not ready to
defend, whatever its length.

\begin{center}
\small
\begin{tabular}{@{}L{11.4cm}ccc@{}}
\toprule
 & \textbf{0} & \textbf{1} & \textbf{2} \\
\midrule
The gap is evidenced by a reproducible search, not asserted & & & \\
Each objective is checkable & & & \\
The design is justified against a named rejected alternative & & & \\
Sample size follows from a power analysis or saturation argument & & & \\
Every measure has validity evidence & & & \\
The analysis is specified per research question, before any data & & & \\
Threats are graded, with directions & & & \\
A null result would still be worth reporting, and you can say why & & & \\
Ethics route and consent are settled & & & \\
The timeline has slack and names its dependencies & & & \\
\bottomrule
\end{tabular}
\end{center}
\end{conceptbox}
